% ============================================================
% DEMO 3 - Elsevier journal template (elsarticle)
% Live demo: show how a .cls changes the entire layout
%
% Compile with:
%   xelatex paper.tex
%   biber paper
%   xelatex paper.tex
%   xelatex paper.tex
% ============================================================

\documentclass[preprint, 12pt, authoryear]{elsarticle}

% ----- Packages -----
\usepackage{amsmath, amssymb}
\usepackage{booktabs}
\usepackage{graphicx}
\usepackage[colorlinks=true, linkcolor=blue, citecolor=blue]{hyperref}

% ----- Custom commands -----
\newcommand{\E}{\mathbb{E}}
\newcommand{\Var}{\mathrm{Var}}

% ============================================================
\begin{document}

\begin{frontmatter}

\title{Survival Analysis with Machine Learning:\\
       A Comparative Study}

\author[unilab]{Imad El Badisy\corref{cor1}}
\ead{elbadisyimad@gmail.com}

\cortext[cor1]{Corresponding author}
\affiliation[unilab]{
  text={Department of Biostatistics,\\
        University of Medicine}
}

\begin{abstract}
We compare survival machine learning methods on a benchmark
dataset. Cox regression serves as baseline. Results show that
ensemble methods improve concordance by 5--8\%.
\end{abstract}

\begin{keyword}
  survival analysis \sep machine learning \sep
  concordance index \sep Cox regression
\end{keyword}

\end{frontmatter}

% ============================================================

\section{Introduction}
\label{sec:intro}

Survival analysis models time-to-event outcomes, which arise
frequently in clinical and epidemiological research
\citep{cox1972}.
Classical methods such as the Cox proportional hazards model
\citep{cox1972} assume a multiplicative hazard structure.
Recent approaches extend this to machine learning
\citep{hastie2009}.

\section{Methods}
\label{sec:methods}

\subsection{Notation}

Let $(T_i, \delta_i, \mathbf{x}_i)$ denote the observed time,
event indicator, and covariates for subject $i = 1,\ldots,n$.

\subsection{Cox Proportional Hazards}

The Cox model specifies:
\begin{equation}
  h(t \mid \mathbf{x})
    = h_0(t) \exp(\mathbf{x}^\top \boldsymbol{\beta})
  \label{eq:cox}
\end{equation}

Parameters are estimated by maximizing the partial likelihood.

\subsection{Concordance Index}

Model discrimination is measured by Harrell's C-index:
\begin{equation}
  C = \frac{\sum_{i:\delta_i=1}
            \sum_{j: T_j > T_i}
            \mathbf{1}(\hat{r}_i > \hat{r}_j)}
           {\sum_{i:\delta_i=1}
            \sum_{j: T_j > T_i} 1}
\end{equation}

\section{Results}
\label{sec:results}

Table~\ref{tab:results} summarizes concordance indices
across all methods.

\begin{table}[htbp]
  \centering
  \caption{Concordance index (C-index) by method.
           Values are mean $\pm$ SD over 10 cross-validation folds.}
  \label{tab:results}
  \begin{tabular}{lcc}
    \toprule
    Method            & C-index         & IBS             \\
    \midrule
    Cox PH            & $0.632 \pm 0.018$ & $0.184 \pm 0.009$ \\
    Random Survival Forest & $0.681 \pm 0.015$ & $0.162 \pm 0.008$ \\
    Gradient Boosting & $0.688 \pm 0.014$ & $0.158 \pm 0.007$ \\
    Deep Survival     & $0.671 \pm 0.017$ & $0.169 \pm 0.009$ \\
    \bottomrule
  \end{tabular}
\end{table}

\section{Discussion}
\label{sec:discussion}

The results in Table~\ref{tab:results} confirm that ensemble
methods outperform Cox regression on this benchmark.
The Cox model \eqref{eq:cox} remains valuable for its
interpretability \citep{hastie2009}.

\section{Conclusion}

Machine learning methods provide meaningful improvements
over Cox regression. Future work should assess calibration.

% ----- Bibliography -----
\bibliographystyle{elsarticle-harv}
\bibliography{../04-bibliography/refs}

\end{document}
